\subsection{Pose Recovery and Mesh Reconstruction}\label{subsec:mesh_recon}

The \texttt{ManoHand} class in \texttt{mano\_mesh\_reconstruction} ties the solver
to the forward model for a single hand. On construction it instantiates a
\texttt{KinematicModel} with the \texttt{MANOArmature}, evaluates it at the zero
pose to obtain the rest-pose keypoints, and reorders those keypoints from
\acrshort{mano} ordering into MediaPipe ordering. This reordered template shares
its indexing with the detector's output, so that template and target are aligned
when both are passed to the solver.

The per-frame work is performed by \texttt{solve}. The observed landmarks are
assembled into a target skeleton and normalised to the template's scale: the ratio
of the wrist-to-middle-finger-base distance in the template to the same distance
in the target is applied to the target, which is then translated so that its wrist
coincides with the template's. The normalised target and the template are passed
to \texttt{adaptive\_IK}, which returns the joint rotations as matrices. Each
matrix is converted to an axis-angle vector, with guards for degenerate cases: a
zero matrix is treated as the identity, and a conversion that yields a
non-numeric result or raises an error falls back to a zero rotation. The resulting
axis-angle pose, whose first entry is the global wrist rotation, is applied to the
forward model to deform the mesh. \texttt{solve} returns both the mesh vertices,
for local visualisation, and the axis-angle pose, which is the quantity
transmitted to the renderer.
